\section{Open Questions}

Five substantive questions remain unresolved after the replication. Each is stated with the basis (what in the paper or in our re-analysis motivates it) and a concrete next-step probe implementable at low cost.

\subsection*{Q1. Radiation-quality generalization}
\textbf{Question.} Does the impaired low-dose p53 activation signature in N2+ generalize to other radiation qualities (proton, alpha, neutron) beyond 0.05/2 Gy X-rays?

\textbf{Basis.} The paper's headline biological finding (HALLMARK\_P53\_PATHWAY fold-enrichment 5.4$\times$ in N0 vs.\ 2.4$\times$ in N2+ at 0.05 Gy) is confined to a single radiation modality. Second-primary-neoplasm survivors typically had heterogeneous prior radiotherapy exposures, so distinguishing dose-driven from LET-driven attenuation is medically relevant.

\textbf{Next steps.} Mine ENA/GEO for other bulk-RNA-seq fibroblast irradiation datasets across LET regimes (NASA space-radiation, GSI heavy-ion series). Re-run the same MSigDB ORA against those DEG lists at matched delivered doses; compare fold-enrichment between control and cancer-survivor cohorts.

\subsection*{Q2. Single-cell resolution vs cell-composition confound}
\textbf{Question.} Is the group-level DEG signal driven by a small number of donors or by altered cell-type composition, or does it hold at single-cell resolution across fibroblast subpopulations?

\textbf{Basis.} Bulk RNA-seq averages over a mix of fibroblast subtypes (dermal papillary vs.\ reticular, myofibroblast contamination, senescent subpopulations). A 5.4$\times$ vs.\ 2.4$\times$ fold-enrichment difference could reflect composition, not per-cell response.

\textbf{Next steps.} Deconvolve bulk profiles with CIBERSORTx or MuSiC against a public single-cell fibroblast atlas (Tabula Sapiens skin fraction). Test whether composition-corrected DE still yields the group difference. If a single-cell irradiation dataset on primary fibroblasts becomes available, re-run at single-cell resolution.

\subsection*{Q3. Tumor-stroma translation}
\textbf{Question.} Are the KiKme DEG signatures diagnostic or prognostic when applied to the tumor stroma of the corresponding second primaries?

\textbf{Basis.} The paper uses skin fibroblasts as a proxy for a systemic radiosensitivity phenotype. Whether the same signature is enriched in stromal compartments of second-primary tumors from these patients (or matched cohorts) is unaddressed.

\textbf{Next steps.} Extract stromal-compartment gene expression from TCGA or the Human Tumor Atlas Network for the second-primary neoplasm types over-represented in KiKme, score against the KiKme low-dose signature (ssGSVA), and correlate with clinical outcome using public follow-up data.

\subsection*{Q4. Normalization / DE-method sensitivity}
\textbf{Question.} How sensitive are the exact DEG counts and top-12 gene list to choice of normalization and DE-model (voom/limma vs.\ DESeq2 vs.\ edgeR-glmQLF)?

\textbf{Basis.} The paper uses voom/limma exclusively. Our replication was blocked from method-sensitivity testing because the raw count matrix is not deposited. If the group-difference signal is method-specific, the biological claim weakens; if method-robust, it strengthens.

\textbf{Next steps.} Request the count matrix under a DUA from Mainz IMBEI. Re-run DE with DESeq2 (Wald) and edgeR-glmQLF. Compute Jaccard overlap of DEG sets vs.\ the published voom/limma set. Specifically test whether the N0 vs.\ N2+ fold-enrichment difference at HALLMARK\_P53\_PATHWAY holds across all three methods.

\subsection*{Q5. Per-donor predictive utility}
\textbf{Question.} Can a machine-learning classifier trained on the KiKme low-dose DEG profile predict second-primary-neoplasm risk from a fibroblast transcriptome alone?

\textbf{Basis.} The paper reports group-level enrichment differences but does not build a per-individual model. Whether the transcriptomic signature carries enough per-donor signal to serve as a clinical radiosensitivity biomarker is unresolved and clinically consequential.

\textbf{Next steps.} Using the 8{,}134-gene AF1 union DEG matrix as a feature space, train a regularized logistic regression (glmnet) or small MLP with donor-level cross-validation to classify N0 vs.\ N2+ from the 0.05 Gy transcriptome. Report AUC, calibration, and per-gene attributions; benchmark against a random-feature baseline to quantify how much of the group difference is per-donor decodable vs.\ group-mean-only.
